Suppose $\ell\mid k_n$. Proposition~\ref{prop:F} gives $\ell^{d_{f_0}}\le C_n$ for the saturation component $f_0$, and Lemma~\ref{lem:B} gives $d_{f_0}=d_n(\ell)$; hence $\ell^{d_n(\ell)}\le C_n$. Contrapositively, $\ell^{d_n(\ell)}>C_n$ implies $\ell\nmid k_n$. The discrete core of this argument (saturation exists, every lattice vector obeys the Blichfeldt bound in the power-$m/2$ form, saturation forces height $\ge L$, hence $C<\ell^d$ is impossible) is the Lean theorem \path{TheoremAInputs.theoremA_of_inputs}.
